\begin{prob}
    Suppose a binary search tree has been augmented so that each node contains an
    additional attribute called \python{size} which contains the number of nodes
    in the subtree rooted at that node. Complete the following code
    so that it computes the value of the $k$th smallest key in the subtree
    rooted at \python{node}, where $k = 1$ is the minimum.

    \inputminted{python}{\thisdir/include/order_statistic.py}

    \begin{soln}
        If \python{order > k}, we should check the subtree of the left child. We
        are looking for the $k$th smallest thing among the nodes in the left subtree,
        so our recursive call is \python{order_statistic(node.left, k)}.

        On the other hand,
        if \python{order < k}, we should check the subtree of the right child, but we
        have to ``update'' the value of $k$. We do not want the $k$th smallest thing
        in the right subtree; rather, we want the \python{k - order} smallest thing (or,
        equivalently, \python{k - left_size - 1}).

        To see this, consider a simple example. Suppose we want the $k = 7$ smallest
        key in the tree and that \python{node.left.size} is 4. This makes the key
        of the current node the fifth smallest in the tree. All keys in the right subtree
        are larger, so we want the second smallest among them.

        The recursive call in this case is therefore \python{order_statistic(node.right, k - order)}.
    \end{soln}
\end{prob}
